% concatenated from main__5_.tex
\documentclass[11pt,a4paper]{article}
\pdfoutput=1
\usepackage[utf8]{inputenc}
\usepackage[T1]{fontenc}
\usepackage{amsmath,amssymb,amsthm}
\usepackage{geometry}
\usepackage{hyperref}
\usepackage{enumitem}

\geometry{margin=2.5cm}
\newtheorem{theorem}{Theorem}[section]
\newtheorem{lemma}[theorem]{Lemma}
\newtheorem{condition}{Condition}

\title{Parametric Statistical Inference in the Zone of Moderate Deviation Probabilities}
\author{M.\,S.~Ermakov}
\date{16 August 2026}

\begin{document}
\maketitle

\begin{abstract}
We study parametric statistical inference in the zone of moderate deviation probabilities. Using the Hellinger-distance technique for the expansion of the log-likelihood ratio, we extend classical results on normal approximation of maximum-likelihood estimators, Bayesian estimators, and posterior concentration to the moderate-deviation regime under substantially weaker conditions than those previously required.
\end{abstract}

\section{Introduction}

Small probability values constantly appear in statistics. Significance levels in confidence estimation and error probabilities in hypothesis testing take small values. The usual approach to the study of such problems consists in treating them as problems concerning probabilities of large and moderate deviations. The question arises: ``under what conditions can results on the normal approximation of the distributions of statistics be extended to the zone of moderate deviation probabilities?''

For maximum-likelihood estimators and likelihood-ratio tests such results have been obtained in a number of papers under rather stringent conditions~\cite{1,2,8,11,14,17} compared with those under which statements on their asymptotic normality were proved~\cite{10,12,13,18}. For Bayesian estimators only concentration inequalities were obtained~\cite{10}. The logarithmic asymptotics corresponding to these inequalities~\cite{10} is optimal in order.

For the study of the asymptotics of moderate-deviation probabilities of maximum-likelihood estimators and likelihood-ratio statistics the Taylor expansion of the log-likelihood ratio based on its derivatives was previously used~\cite{1,2,8,11,14,17}. In the proof of local asymptotic normality of the likelihood ratio, along with this technique one also employs the Taylor expansion of the log-likelihood ratio based on the Hellinger distance~\cite{9,10,13,18}. We use the latter technique for the study of statistical inference problems in the zone of moderate deviation probabilities. We succeed in extending to the moderate-deviation zone the principal results obtained by this approach for the normal approximation. We prove a theorem on the local uniform normal approximation of the log-likelihood ratio in the moderate-deviation zone, and we prove large-deviation principles in the moderate-deviation zone for maximum-likelihood estimators and Bayesian estimators. We also prove a large-deviation principle in the moderate-deviation zone for the concentration of the posterior Bayesian measure. The conditions of the theorems turn out to be much weaker than the conditions under which theorems on moderate-deviation probabilities of maximum-likelihood estimators and likelihood-ratio statistics were previously proved~\cite{1,2,8,11,14,17}.

The results mentioned above are stated in Section~\ref{sec:results}. Their proofs are given in the subsequent sections.

We use the letters $c$, $C$ and $c$, $C$ with indices to denote positive constants. We write $\mathbf{1}(A)$ for the indicator of an event $A$. For any event $A$ we denote by $\overline{A}$ its complement. For any two sequences of positive numbers $a_n$ and $b_n$, the relation $a_n=o(b_n)$ means that $a_n/b_n\to 0$ as $n\to\infty$. For any vector $u$ we denote by $u^T$ the transposed vector and by $|u|$ the Euclidean length of $u$.

\section{Problem Setting}
\label{sec:setting}

Let $X_1,\dots,X_n$ be a sample of independent identically distributed random variables having probability measure $P_\theta$, $\theta\in\Theta$, defined on a $\sigma$-algebra $\mathcal{B}$ of a set $S$. The set $\Theta$ is an open bounded set in $\mathbb{R}^d$. The value of the parameter $\theta$ is unknown.

Assume that the probability measures $P_\theta$, $\theta\in\Theta$, have a common support equal to the set $S$ and are absolutely continuous with respect to a measure $\nu$ defined on the same $\sigma$-algebra $\mathcal{B}$. Denote
\begin{equation}
f(x,\theta)=\frac{dP_\theta}{d\nu}(x),\qquad x\in S,
\label{eq:density}
\end{equation}
the density of the probability measure $P_\theta$ with respect to $\nu$.

We shall be interested in statistical inference about a parameter $\theta_0\in\Theta$.

For all $\theta_0$, $\theta_0+\tau\in\Theta$ define the function
\begin{equation}
g(x,\tau)=g(x,\theta_0,\theta_0+\tau)=\Bigl(\frac{f(x,\theta_0+\tau)}{f(x,\theta_0)}\Bigr)^{1/2}-1,
\label{eq:g}
\end{equation}
where $x\in S$.

We say that the statistical experiment $\mathcal{E}=\{(S,\mathcal{B}),P_\theta,\theta\in\Theta\}$ has finite Fisher information at the point $\theta_0\in\Theta$ if there exist a vector-valued function $\phi=\phi_{\theta_0}:S\to\mathbb{R}^d$ and an increasing function $\omega(|\tau|)=\omega_{\theta_0}(|\tau|)\to 0$ as $\tau\to 0$ such that
\begin{equation}
\int_S\bigl(g(x,\tau)-\tau^T\phi(x)\bigr)^2\,dP_{\theta_0}<|\tau|^2\omega(|\tau|).
\label{eq:FI}
\end{equation}
Define the Fisher information matrix
\begin{equation}
I(\theta_0)=4\int_S\phi\phi^T\,dP_{\theta_0}.
\label{eq:I}
\end{equation}
We shall assume that the rank of the matrix $I(\theta)$ is equal to $d$ at every point $\theta\in\Theta\subset\mathbb{R}^d$.

We introduce the following condition, which is essentially a condition of identifiability of the parameter value (see condition~(3.2), Chapter~3,~\cite{10}).

\begin{condition}[A0]
For any compact $K\subset\Theta$ and any $\delta>0$,
\begin{equation}
\inf_{\theta\in K}\inf_{\substack{\tau:\theta+\tau\in\Theta\\|\tau|\ge\delta}}
\int_S\bigl(f^{1/2}(x,\theta+\tau)-f^{1/2}(x,\theta)\bigr)^2\,d\nu>0.
\label{eq:A0}
\end{equation}
\end{condition}

Let a sequence $u_n>0$ be given with $u_n\to 0$ and $nu_n^2\to\infty$ as $n\to\infty$. We shall be interested in the moderate-deviation probabilities
\[
P_{\theta_0}\bigl(|\hat\theta_n-\theta_0|>u_n\bigr),
\]
where $\hat\theta_n$ is a maximum-likelihood estimator or a Bayesian estimator. We shall also study the moderate-deviation probabilities of the concentration of the posterior Bayesian measure on the region $|\theta-\theta_0|>u_n$.

The following conditions A1 and A2 (respectively B1 and B2) yield the statements of the theorems according to the rate at which the function $\omega(u)$ tends to zero as $u\to 0$. The weaker moment conditions A1 (respectively B1) will be assumed when $\omega(u)<Cu^\lambda$ for some $\lambda>0$.

\begin{condition}[A1]
For any $\theta_0\in\Theta$, for some neighbourhood $U\subset\Theta$ of the point $\theta_0$, there exist positive constants $C$, $\gamma$, a sequence $\kappa_n=\kappa_n(u_n)$ with $\kappa_n\to 0$ as $n\to\infty$, and a sequence of functions $h_n:S\to\mathbb{R}$ such that for any $C_1$ and for $n>n_0(C_1)$ one has
\begin{equation}
\sup_{|\tau|<C_1u_n}\Bigl|\log\frac{f(x,\theta+\tau)}{f(x,\theta)}\Bigr|\le\kappa_nh_n(x),\qquad x\in S,
\label{eq:A1}
\end{equation}
and $\mathbb{E}_\theta\exp\{\gamma h_n(X_1)\}<C$ for all $\theta\in U$.
\end{condition}

\begin{condition}[A2]
For any $\theta_0\in\Theta$, for some neighbourhood $U\subset\Theta$ of the point $\theta_0$, there exist positive constants $C$, $c$, $\gamma$, a sequence $\kappa_n=\kappa_n(u_n)$ with $\kappa_n|\log u_n|<c$, and a sequence of functions $h_n:S\to\mathbb{R}$ such that~\eqref{eq:A1} holds and $\mathbb{E}_\theta\exp\{\gamma h_n(X_1)\}<C$ for all $\theta\in U$.
\end{condition}

Conditions A1 and A2 with $\kappa_n=\gamma_n^{-1}$ are sufficient for the more difficult-to-verify conditions B1 and B2, respectively, under which the statements of the theorems will be proved.

\begin{condition}[B1]
For any $\theta_0\in\Theta$, for some neighbourhood $U\subset\Theta$ of the point $\theta_0$, there exist constants $\varepsilon>0$, $C$ and a sequence $\gamma_n=\gamma_n(u_n)$ with $\gamma_n\to\infty$ as $n\to\infty$ such that for any $C_1$ and for $n>n_0(C_1)$ one has
\begin{equation}
\mathbb{E}_\theta\Biggl(\frac{f(X,\theta+\tau)}{f(X,\theta)}\,\mathbf{1}\Bigl(\Bigl|\log\frac{f(X,\theta+\tau)}{f(X,\theta)}\Bigr|>\varepsilon\Bigr)\Biggr)^{\gamma_n}<C
\label{eq:B1}
\end{equation}
for all $\theta\in U$, $|\tau|<C_1u_n$.
\end{condition}

\begin{condition}[B2]
For any $\theta_0\in\Theta$, for some neighbourhood $U\subset\Theta$ of the point $\theta_0$, there exist constants $\varepsilon>0$, $c$, $C$ and a sequence $\gamma_n=\gamma_n(u_n)$ with $\gamma_n>c|\log u_n|$ such that~\eqref{eq:B1} holds for all $\theta\in U$, $|\tau|<C_1u_n$.
\end{condition}

In the proofs of the theorems conditions A1 and A2 are used when
\[
\Bigl|\log\frac{f(X,\theta+\tau)}{f(X,\theta)}\Bigr|>\varepsilon.
\]
This means that it is in fact sufficient to require their validity on the sets $\{h_n(x)>\varepsilon/\kappa_n\}$.

For the proof of the theorems under condition A1 or B1 we need in addition the fulfilment of conditions C1 and C2.

\begin{condition}[C1]
For any $\theta_0\in\Theta$ there exist constants $\gamma>0$, $\lambda>0$ and $C$ such that for all $u$, $0<u<u_0$, one has $\omega(u)<Cu^\lambda$ and for all $\theta$ from some neighbourhood $U$ of the point $\theta_0$ one has
\begin{equation}
\mathbb{E}_\theta\bigl[\phi^2(X_1)\mathbf{1}(|\phi(X_1)|>\varepsilon u^{-1})\bigr]<Cu^\gamma.
\label{eq:C1}
\end{equation}
\end{condition}

\begin{condition}[C2]
For any $\theta_0\in\Theta$ there exist $\gamma>0$ and $C>0$ such that for all $\tau\in\mathbb{R}^d$, $0<|\tau|<u_0$, $\theta_0+\tau\in\Theta$, one has
\begin{equation}
\bigl|\mathbb{E}_{\theta_0}[g^2(X_1,\tau)]-\tfrac14\tau^T I(\theta_0)\tau\bigr|<C|\tau|^{2+\gamma}.
\label{eq:C2}
\end{equation}
\end{condition}

Until the present work, proofs of theorems on moderate-deviation probabilities of maximum-likelihood estimators and likelihood-ratio statistics relied on the Taylor expansion of the log-likelihood ratio
\begin{equation}
\log\frac{f(X,\theta+u)}{f(X,\theta)}=u\frac{\partial\log f(X,\theta+u)}{\partial u}\Big|_{u=0}-\frac{u^2}{2}\frac{\partial^2\log f(X,\theta+u)}{\partial u^2}\Big|_{u=v}
\label{eq:Taylor}
\end{equation}
with $v\in(0,u)$. Assumptions were made about the existence of exponential moments of the suprema, with respect to the parameter, of the derivatives of the log-likelihood ratio. Such conditions were considered for derivatives up to the third order~\cite{1,2,8,11,14,17}.

The weakest assumptions appear to be contained in~\cite{1}. In~\cite{1} a theorem on the moderate-deviation probability of the maximum-likelihood estimator was proved under the assumption of existence of an exponential moment of the supremum, with respect to the parameter, of the second derivative of the log-likelihood ratio. The same statement was proved in~\cite{1} when only the existence of an exponential moment of the supremum of the first derivative was required, but it was assumed that the likelihood-ratio function of the sample has, with probability one, a unique local maximum coinciding with the global maximum. In the setting of the present paper the condition of existence of an exponential moment of the supremum of the first derivative can be viewed as a particular case of condition A1 with $\kappa_n=u_n$ under some additional mild restrictions. Moreover, we introduce no assumptions about uniqueness of a local maximum of the likelihood-ratio function coinciding with the global maximum.

The following condition is needed when studying maximum-likelihood estimators in the multidimensional case. For Bayesian estimators it is absent.

\begin{condition}[D]
The set $\Theta$ is convex and bounded, and there exists $m>d$ such that
\begin{equation}
\sup_{\theta\in\Theta}\int_S\Bigl(\frac{|\nabla f(x,\theta)|}{f(x,\theta)}\Bigr)^m f(x,\theta)\,d\nu<\infty.
\label{eq:D}
\end{equation}
Here $\nabla f(x,\theta)$ denotes the gradient of the function $f(x,\theta)$ with respect to the parameter $\theta$.
\end{condition}

For $\varepsilon>0$ put $\phi_\varepsilon(x)=\phi(x)\mathbf{1}(u_n|\phi(x)|<\varepsilon)$, $x\in S$.
The following condition will be used when studying the uniform convergence of the log-likelihood ratio to its locally asymptotically normal approximation in the moderate-deviation zone.

\begin{condition}[E]
For any $\theta_0\in\Theta$ there exist $\beta_1>d$ and $\beta_2>d$ such that for some neighbourhood $U$ of the point $\theta_0$ and for any $\theta$, $\theta_0+u$, $\theta_0+v\in U$ one has
\begin{equation}
\mathbb{E}_\theta\Bigl(\log\frac{f(X_1,\theta_0+v)}{f(X_1,\theta_0+u)}-2(v-u)^T\phi_\varepsilon(X_1)\Bigr)^{\beta_1}\le C|v-u|^{\beta_2}.
\label{eq:E}
\end{equation}
\end{condition}

\section{Main Results}
\label{sec:results}

\subsection{Local asymptotic normality of the log-likelihood ratio in the moderate-deviation zone}

For $u\in\mathbb{R}^d$ denote
\begin{equation}
\zeta_n(u)=2u^T\sum_{i=1}^n\phi_\varepsilon(X_i)-\frac12 nu^T I(\theta_0)u
\end{equation}
and
\begin{equation}
\psi_n=n^{-1/2}I^{-1/2}(\theta_0)\sum_{i=1}^n\phi_\varepsilon(X_i),
\end{equation}
where $\phi_\varepsilon(x)=\phi(x)\mathbf{1}(|\phi(x)|<u_n^{-1}\varepsilon)$, $x\in S$.

For $\theta_0$, $\theta_0+u_nb$, $\theta_0+u_nb+u\in\Theta$, $b\in\mathbb{R}^d$, put
\begin{equation}
\xi_n(X_i,u_nb,u)=\log\frac{f(X_i,\theta_0+u_nb+u)}{f(X_i,\theta_0+u_nb)},\qquad 1\le i\le n.
\end{equation}
If $b=0$ we write $\xi_n(X_i,u)=\xi_n(X_i,0,u)$.

The following Theorem~\ref{th:LAN} is an analogue of the theorem on local asymptotic normality of the log-likelihood ratio~\cite{9,10,13,18} for the zone of moderate deviation probabilities.

\begin{theorem}\label{th:LAN}
Let the statistical experiment have finite Fisher information at the point $\theta_0\in\Theta\subset\mathbb{R}^d$ and let conditions A0 and B2 be satisfied. Then for any $C$ there exist positive functions $\omega_1(u)$, $\omega_1(u)\to 0$ as $u\to 0$, and $\omega_2(u)$, $\omega_2(u)\to 0$ as $u\to 0$, such that for $|u|<C u_n$ and $b\in\mathbb{R}^d$ one has
\begin{equation}
P_{\theta_0+u_nb}\Biggl(\Biggl|\sum_{i=1}^n\xi_n(X_i,u_nb,u)-\zeta_n(u)\Biggr|>nu_n^2\omega_1(u_n)\Biggr)<C_1\exp\bigl\{-nu_n^2/\omega_2(u_n)\bigr\}.
\label{eq:13}
\end{equation}
If in addition conditions D and E are satisfied, then for any $C$ there exist positive functions $\omega_3(u)$, $\omega_3(u)\to 0$ as $u\to 0$, and $\omega_4(u)$, $\omega_4(u)\to 0$ as $u\to 0$, such that
\begin{equation}
P_{\theta_0+u_nb}\Biggl(\sup_{|u|<C u_n}\Biggl|\sum_{i=1}^n\xi_n(X_i,u_nb,u)-\zeta_n(u)\Biggr|>nu_n^2\omega_3(u_n)\Biggr)<C_1\exp\bigl\{-nu_n^2/\omega_4(u_n)\bigr\}.
\label{eq:14}
\end{equation}
All the inequalities above remain valid if condition B2 is replaced by conditions B1, C1 and C2.
\end{theorem}

The reasoning will be based on the following theorem.

\begin{theorem}\label{th:LD}
Let the statistical experiment have finite Fisher information at the point $\theta_0\in\Theta\subset\mathbb{R}^d$. Then for any open set $\Omega\subset\mathbb{R}^d$ one has
\begin{equation}
\lim_{n\to\infty}(nu_n^2)^{-1}\log P_{\theta_0}(\psi_n\in n^{1/2}u_n\Omega)=-\frac12\inf_{u\in\Omega}|u|^2
\label{eq:15}
\end{equation}
and for any vector $b\in\mathbb{R}^d$ one has
\begin{equation}
\lim_{n\to\infty}(nu_n^2)^{-1}\log P_{\theta_0+n^{1/2}u_nb}(\psi_n\in n^{1/2}u_nb+n^{1/2}u_n\Omega)=-\frac12\inf_{u\in\Omega}|u|^2.
\label{eq:16}
\end{equation}
\end{theorem}

The proofs of Theorem~\ref{th:LD} can be found in~\cite{6}, Theorem~3, and in~\cite{7}.

\subsection{Maximum-likelihood estimator}

Henceforth we assume that for all $x\in S$ the density $f(x,\theta)$ is continuous in $\theta$ on the set $\Theta$.

Define the maximum-likelihood estimator $\hat\theta_n=\hat\theta_n(X_1,\dots,X_n)$ by the equation
\begin{equation}
\prod_{i=1}^n f(X_i,\hat\theta_n)=\max_{\theta\in\Theta}\prod_{i=1}^n f(X_i,\theta),
\label{eq:MLE}
\end{equation}
where as the maximum-likelihood estimator $\hat\theta_n$ one takes any of the points at which the maximum is attained.

\begin{theorem}\label{th:MLE}
Let the statistical experiment have finite Fisher information at the point $\theta_0\in\Theta$ and let conditions A0 and B2 be satisfied. Then, if $d=1$ and the set $\Theta$ is a bounded interval in $\mathbb{R}$, for any open set $\Omega\subset\mathbb{R}$ one has
\begin{equation}
\lim_{n\to\infty}(nu_n^2/2)^{-1}\log P_{\theta_0}\bigl(I^{1/2}(\theta_0)(\hat\theta_n-\theta_0)\in u_n\Omega\bigr)=-\inf_{x\in\Omega}|x|^2
\label{eq:18}
\end{equation}
and
\begin{equation}
\lim_{n\to\infty}(nu_n^2/2)^{-1}\log P_{\theta_0+u_n}\bigl(I^{1/2}(\theta_0)(\hat\theta_n-\theta_0-u_n)\in u_n\Omega\bigr)=-\inf_{x\in\Omega}|x|^2.
\label{eq:19}
\end{equation}
If $d>1$ and in addition condition D is satisfied, then for any open set $\Omega\subset\mathbb{R}^d$ relation~\eqref{eq:18} holds, and for any vector $b\in\mathbb{R}^d$ one has
\begin{equation}
\lim_{n\to\infty}(nu_n^2/2)^{-1}\log P_{\theta_0+u_nb}\bigl(I^{1/2}(\theta_0)(\hat\theta_n-\theta_0-u_nb)\in u_n\Omega\bigr)=-\inf_{x\in\Omega}|x|^2.
\label{eq:20}
\end{equation}
For any $\delta>0$ and any vector $b\in\mathbb{R}^d$ one has
\begin{equation}
\lim_{n\to\infty}(nu_n^2/2)^{-1}\log P_{\theta_0+u_nb}\bigl(|I^{1/2}(\theta_0)(\hat\theta_n-\theta_0)-2n^{-1/2}\psi_n|>\delta u_n\bigr)=-\infty,
\label{eq:21}
\end{equation}
\begin{equation}
\begin{aligned}
&\lim_{n\to\infty}(nu_n^2/2)^{-1}\log P_{\theta_0+u_nb}\Biggl(
\Biggl|n^{-1}\sum_{i=1}^n\xi_n(X_i,u_nb,\hat\theta_n-\theta_0-u_nb)\\
&\qquad-2(\psi_n-n^{1/2}u_nb)^T(\psi_n-n^{1/2}u_nb)\Biggr|>\delta nu_n^2\Biggr)=-\infty
\end{aligned}
\label{eq:22}
\end{equation}
and
\begin{equation}
\begin{aligned}
&\lim_{n\to\infty}(nu_n^2/2)^{-1}\log P_{\theta_0+u_nb}\Biggl(
\Biggl|\sum_{i=1}^n\xi_n(X_i,u_nb,\hat\theta_n-\theta_0-u_nb)\\
&\qquad-n(\hat\theta_n-\theta_0-u_nb)^T I(\theta_0)(\hat\theta_n-\theta_0-u_nb)/2\Biggr|>\delta nu_n^2\Biggr)=-\infty.
\end{aligned}
\label{eq:23}
\end{equation}
All the statements above remain valid if condition B2 is replaced by conditions B1, C1 and C2.
\end{theorem}

The following theorem is a consequence of Theorem~\ref{th:MLE}. It treats the large-deviation probabilities of the maximum-likelihood estimator in the setting of local asymptotic efficiency in the Bahadur sense.

\begin{theorem}\label{th:MLE-Bahadur}
Let the statistical experiment have finite Fisher information at the point $\theta_0\in\Theta$ and let conditions A0 and B1 be satisfied. Then, if $d=1$ and the set $\Theta$ is a bounded interval in $\mathbb{R}$, for any open set $\Omega\subset\mathbb{R}$ one has
\begin{equation}
\lim_{u\to 0}\lim_{n\to\infty}(nu^2/2)^{-1}\log P_{\theta_0}\bigl(I^{1/2}(\theta_0)(\hat\theta_n-\theta_0)\in u\Omega\bigr)=-\inf_{x\in\Omega}|x|^2.
\label{eq:24}
\end{equation}
If $d>1$ and in addition condition D is satisfied, then for any open set $\Omega\subset\mathbb{R}^d$ relation~\eqref{eq:24} holds.

Moreover, for any $\delta>0$ one has
\begin{equation}
\lim_{u\to 0}\lim_{n\to\infty}(nu^2/2)^{-1}\log P_{\theta_0}\bigl(|I^{1/2}(\theta_0)(\hat\theta_n-\theta_0)-2n^{-1/2}\psi_n|>\delta u\bigr)=-\infty,
\label{eq:25}
\end{equation}
\begin{equation}
\lim_{u\to 0}\lim_{n\to\infty}(nu^2/2)^{-1}\log P_{\theta_0}\Biggl(
\Biggl|n^{-1}\sum_{i=1}^n\xi_n(X_i,\hat\theta_n-\theta_0)-2\psi_n^T\psi_n\Biggr|>\delta nu^2\Biggr)=-\infty
\label{eq:26}
\end{equation}
and
\begin{equation}
\lim_{u\to 0}\lim_{n\to\infty}(nu^2/2)^{-1}\log P_{\theta_0}\Biggl(
\Biggl|\sum_{i=1}^n\xi_n(X_i,\hat\theta_n-\theta_0)
-n(\hat\theta_n-\theta_0)^T I(\theta_0)(\hat\theta_n-\theta_0)/2\Biggr|>\delta nu^2\Biggr)=-\infty.
\label{eq:27}
\end{equation}
\end{theorem}

For the moderate-deviation problem the relations~\eqref{eq:21}--\eqref{eq:23} and~\eqref{eq:25}--\eqref{eq:27} show the equivalence of the likelihood-ratio, Wald and Rao criteria for multidimensional testing of hypotheses about the value of a parameter.

In Theorems~\ref{th:MLE} and~\ref{th:MLE-Bahadur} the set $\Theta$ is assumed bounded. In Theorems~5.1 and~5.4 of Chapter~1 of~\cite{10} inequalities for moderate-deviation probabilities of maximum-likelihood estimators were obtained when the parameter set is unbounded. These inequalities allow one to extend the proofs of Theorems~\ref{th:MLE} and~\ref{th:MLE-Bahadur} to the case of an unbounded set $\Theta$. Thus relations~\eqref{eq:18}--\eqref{eq:27} remain valid for unbounded domains $\Theta$ if in addition conditions~(1)--(3) of Theorem~5.4 of Chapter~1 of~\cite{10} are satisfied. It is easy to see that they can be replaced by the following simpler sufficient conditions.

There exists $\gamma>0$ such that for any compact $K\subset\Theta$ there is $c=c(K)$ with
\begin{equation}
\sup_{|h|>R}\int_S f^{1/2}(x,\theta)f^{1/2}(x,\theta+h)\,d\nu\le c R^{-\gamma},\qquad\theta\in K.
\label{eq:28}
\end{equation}
In addition, for some $m>0$ the following conditions must hold:
\begin{equation}
\begin{gathered}
\sup_{|\theta|<R}\mathbb{E}_\theta|\phi_\theta(X_1)|^2<C(1+R^m),\\
\inf_{\theta\in\Theta}\mathbb{E}_\theta\phi_\theta^2(X_1)>c>0,\\
\sup_{|\theta|<R}\sup_h\omega_\theta(h)<C(1+R^m).
\end{gathered}
\label{eq:29}
\end{equation}

Theorem~5.4 of Chapter~1 of~\cite{10} is stated and proved for $\Theta\subset\mathbb{R}$. However it remains valid for $\Theta\subset\mathbb{R}^d$ if condition D is assumed in addition. In the multidimensional case the proofs differ only in the estimate of the modulus of continuity of the likelihood ratio, which is provided by condition D.

\subsection{Bayesian estimator}

Let $l_1:\mathbb{R}_+\to\mathbb{R}_+$ be an increasing function such that $l_1(0)=0$ and $l_1(x)>0$ if $x\neq 0$. Suppose that $l_1$ satisfies the following condition: there exist positive constants $C_1$, $C_2$ and $\kappa_1$, $\kappa_2$ such that for any $a>1$ and $x>0$
\begin{equation}
l_1(ax)\le C_1 a^{\kappa_1}l_1(x),\qquad
C_2(a-1)^{\kappa_2}<\frac{l_1(ax)-l_1(x)}{l_1(x)}.
\label{eq:30}
\end{equation}
The loss function $l(u)$ has the form $l(u)=l_1(\|u\|)$, where $\|u\|$ is some (not necessarily Euclidean) norm of the vector $u\in\mathbb{R}^d$.

We assume that the Bayesian prior density $\pi(\theta)$, $\theta\in\Theta$, is continuous, bounded and $\pi(\theta)>0$ for all $\theta\in\Theta$.

The Bayesian estimator of the parameter $\theta$ is defined by
\begin{equation}
\tilde\theta_n=\arg\min_{t\in\Theta}\int_\Theta l(t-\theta)\,q_n(\theta)\,d\theta,
\label{eq:Bayes}
\end{equation}
where
\begin{equation}
q_n(\theta)=\frac{\prod_{j=1}^n f(X_j,\theta)\,\pi(\theta)}{\int_\Theta\prod_{j=1}^n f(X_j,\theta)\,\pi(\theta)\,d\theta}
\label{eq:posterior}
\end{equation}
is the posterior density.

The statements of the following Theorems~\ref{th:Bayes} and~\ref{th:Bayes-Bahadur} differ from those of Theorems~\ref{th:MLE} and~\ref{th:MLE-Bahadur} essentially only by the absence of condition D and by the replacement of the maximum-likelihood estimator $\hat\theta_n$ by the Bayesian estimator $\tilde\theta_n$.

\begin{theorem}\label{th:Bayes}
Let the statistical experiment have finite Fisher information at the point $\theta_0\in\Theta\subset\mathbb{R}^d$ and let conditions A0 and B2 be satisfied. Then for any open set $\Omega\subset\mathbb{R}^d$ and any vector $b\in\mathbb{R}^d$ one has
\begin{equation}
\lim_{n\to\infty}(nu_n^2/2)^{-1}\log P_{\theta_0+u_nb}\bigl(I(\theta_0)^{1/2}(\tilde\theta_n-\theta_0-u_nb)\in u_n\Omega\bigr)=-\inf_{x\in\Omega}|x|^2.
\label{eq:33}
\end{equation}
For any $\delta>0$ and any vector $b\in\mathbb{R}^d$ one has
\begin{equation}
\lim_{n\to\infty}(nu_n^2/2)^{-1}\log P_{\theta_0+u_nb}\bigl(|I^{1/2}(\theta_0)(\tilde\theta_n-\theta_0)-2n^{-1/2}\psi_n|>\delta u_n\bigr)=-\infty,
\label{eq:34}
\end{equation}
\begin{equation}
\begin{aligned}
&\lim_{n\to\infty}(nu_n^2/2)^{-1}\log P_{\theta_0+u_nb}\Biggl(
\Biggl|n^{-1}\sum_{i=1}^n\xi_n(X_i,u_nb,\tilde\theta_n-\theta_0-u_nb)\\
&\qquad-2(\psi_n-n^{1/2}u_nb)^T(\psi_n-n^{1/2}u_nb)\Biggr|>\delta nu_n^2\Biggr)=-\infty
\end{aligned}
\label{eq:35}
\end{equation}
and
\begin{equation}
\begin{aligned}
&\lim_{n\to\infty}(nu_n^2/2)^{-1}\log P_{\theta_0+u_nb}\Biggl(
\Biggl|\sum_{i=1}^n\xi_n(X_i,u_nb,\tilde\theta_n-\theta_0-u_nb)\\
&\qquad-n(\tilde\theta_n-\theta_0-u_nb)^T I(\theta_0)(\tilde\theta_n-\theta_0-u_nb)/2\Biggr|>\delta nu_n^2\Biggr)=-\infty.
\end{aligned}
\label{eq:36}
\end{equation}
All the statements above remain valid if condition B2 is replaced by conditions B1, C1 and C2.
\end{theorem}

\begin{theorem}\label{th:Bayes-Bahadur}
Let the statistical experiment have finite Fisher information at the point $\theta_0\in\Theta\subset\mathbb{R}^d$ and let conditions A0 and B1 be satisfied. Then for any open set $\Omega\subset\mathbb{R}^d$ one has
\begin{equation}
\lim_{u\to 0}\lim_{n\to\infty}(nu^2/2)^{-1}\log P_{\theta_0}\bigl(I(\theta_0)^{1/2}(\tilde\theta_n-\theta_0)\in u\Omega\bigr)=-\inf_{x\in\Omega}|x|^2.
\label{eq:37}
\end{equation}
For any $\delta>0$ one has
\begin{equation}
\lim_{u\to 0}\lim_{n\to\infty}(nu^2/2)^{-1}\log P_{\theta_0}\bigl(|I^{1/2}(\theta_0)(\tilde\theta_n-\theta_0)-2n^{-1/2}\psi_n|>\delta u\bigr)=-\infty,
\label{eq:38}
\end{equation}
\begin{equation}
\lim_{u\to 0}\lim_{n\to\infty}(nu^2/2)^{-1}\log P_{\theta_0}\Biggl(
\Biggl|n^{-1}\sum_{i=1}^n\xi_n(X_i,\tilde\theta_n-\theta_0)-2\psi_n^T\psi_n\Biggr|>\delta nu^2\Biggr)=-\infty
\label{eq:39}
\end{equation}
and
\begin{equation}
\lim_{u\to 0}\lim_{n\to\infty}(nu^2/2)^{-1}\log P_{\theta_0}\Biggl(
\Biggl|\sum_{i=1}^n\xi_n(X_i,\tilde\theta_n-\theta_0)
-n(\tilde\theta_n-\theta_0)^T I(\theta_0)(\tilde\theta_n-\theta_0)/2\Biggr|>\delta nu^2\Biggr)=-\infty.
\label{eq:40}
\end{equation}
\end{theorem}

The inequalities of Theorems~5.2 and~5.6 of Chapter~1 of~\cite{10} allow one to extend Theorems~\ref{th:Bayes} and~\ref{th:Bayes-Bahadur} to the case when the set $\Theta\subset\mathbb{R}^d$ is unbounded. Theorems~\ref{th:Bayes} and~\ref{th:Bayes-Bahadur} remain valid for an unbounded set $\Theta\subset\mathbb{R}^d$ if in addition conditions~(1)--(3) of Theorem~5.6 of Chapter~1 of~\cite{10} are satisfied. Sufficient conditions for their fulfilment are the conditions~\eqref{eq:28} and~\eqref{eq:29}.

Theorem~5.6 of Chapter~1 of~\cite{10} is proved for the case $l(u)=l_1(|u|)$. However the proof extends essentially without changes to the case $l(u)=l_1(\|u\|)$. We omit the corresponding arguments.

\subsection{On moderate-deviation probabilities of the concentration of the posterior Bayesian measure}

For any measurable set $A\subset\Theta$ the posterior Bayesian measure of the set $A$ is equal to
\begin{equation}
Q_n(A)=\int_A q_n(\theta)\,d\theta.
\label{eq:Qn}
\end{equation}

\begin{theorem}\label{th:posterior}
Let the statistical experiment have finite Fisher information at the point $\theta_0\in\Theta\subset\mathbb{R}^d$ and let conditions A0 and B2 be satisfied. Then for any open set $\Omega\subset\mathbb{R}^d$, for any constant $c$, $0<c<1$, and any vector $b\in\mathbb{R}^d$ one has
\begin{equation}
\lim_{n\to\infty}(nu_n^2/2)^{-1}\log P_{\theta_0+u_nb}\bigl(Q_n(\theta_0+u_n I(\theta_0)^{-1/2}\Omega-u_nb)>c\bigr)=-\inf_{x\in\Omega}|x|^2.
\label{eq:42}
\end{equation}
Relation~\eqref{eq:42} remains valid if condition B2 is replaced by conditions B1, C1 and C2.
\end{theorem}

The following Theorem~\ref{th:posterior-Bahadur} is a consequence of Theorem~\ref{th:posterior} for the setting of asymptotic efficiency in the Bahadur sense.

\begin{theorem}\label{th:posterior-Bahadur}
Let the statistical experiment have finite Fisher information at the point $\theta_0\in\Theta\subset\mathbb{R}^d$ and let conditions A0 and B1 be satisfied. Then for any open set $\Omega\subset\mathbb{R}^d$ and any constant $c$, $0<c<1$, one has
\begin{equation}
\lim_{u\to 0}\lim_{n\to\infty}(nu^2/2)^{-1}\log P_{\theta_0}\bigl(Q_n(\theta_0+u I(\theta_0)^{-1/2}\Omega)>c\bigr)=-\inf_{x\in\Omega}|x|^2.
\label{eq:43}
\end{equation}
\end{theorem}

\section{Proof of Theorems~\ref{th:LAN} and~\ref{th:MLE}}
\label{sec:proofs}

In Subsection~\ref{sub:13} we prove the estimate~\eqref{eq:13} for the approximation of finite-dimensional distributions of the log-likelihood ratio. Then in Subsection~\ref{sub:uniform} we obtain estimates of the uniform approximation, which in particular complete the proof of Theorem~\ref{th:LAN}. In Subsection~\ref{sub:MLE} the proof of Theorem~\ref{th:MLE} will be given.

The proofs of all theorems will be carried out when the true parameter is $\theta_0$, not a sequence of parameters $\theta_0+u_nb$. With this course of reasoning the proofs remain valid for the sequence of parameters $\theta_0+u_nb$ if we replace the function $\phi_{\theta_0}$ by the function $\phi_{\theta_0+u_nb}$. The statement of the lemma below shows that if in the proofs for the parameters $\theta_0+u_nb$ we leave the function $\phi_{\theta_0}$ instead of $\phi_{\theta_0+u_nb}$, then they retain their force.

\begin{lemma}\label{lem:4.1}
Let the statistical experiment have finite Fisher information. Then for $0<|v|<C|u|$ one has
\begin{equation}
\mathbb{E}_{\theta_0+u}\bigl(g(X_1,u,u+v)-v^T\phi_\varepsilon(X_1)\bigr)^2<C|u|^2\omega(C_1|u|).
\end{equation}
\end{lemma}

\begin{proof}
Note that
\[
g(X_1,u,u+v)=\frac{g(X_1,u+v)+1}{g(X_1,u)+1}-1.
\]
Hence
\begin{align*}
&\mathbb{E}_{\theta_0+u}\bigl(g(X_1,u,u+v)-v^T\phi_\varepsilon(X_1)\bigr)^2\\
&=\mathbb{E}_{\theta_0}\Biggl[\Biggl(\frac{g(X_1,u+v)+1}{g(X_1,u)+1}-1-v^T\phi_\varepsilon(X_1)\Biggr)^2(1+g(X_1,u))^2\Biggr]\\
&=\mathbb{E}_{\theta_0}\bigl(g(X_1,u+v)+1-g(X_1,u)-1-v^T\phi_\varepsilon(X_1)(1+g(X_1,u))\bigr)^2\\
&\le C\mathbb{E}_{\theta_0}\bigl(g(X_1,u+v)-(u+v)^T\phi_\varepsilon(X_1)\bigr)^2+C\mathbb{E}_{\theta_0}\bigl(g(X_1,u)-u^T\phi_\varepsilon(X_1)\bigr)^2\\
&\qquad+C|v|^4\mathbb{E}_{\theta_0}|\phi_\varepsilon|^4+C\varepsilon^2\mathbb{E}_{\theta_0}\bigl(g(X_1,u)-u^T\phi_\varepsilon(X_1)\bigr)^2\\
&<C|u|^2\omega(C_1|u|).
\end{align*}
\end{proof}

\subsection{Proof of~\eqref{eq:13}}
\label{sub:13}

We begin with the proof of~\eqref{eq:13} under condition B2. Define the random variables
\[
\xi_{n\varepsilon}(X_i,u)=\xi_n(X_i,u)\mathbf{1}(|\xi_n(X_i,u)|<\varepsilon),\qquad
\bar\xi_{n\varepsilon}(X_i,u)=\xi_n(X_i,u)\mathbf{1}(|\xi_n(X_i,u)|>\varepsilon),\qquad 1\le i\le n.
\]
We have
\begin{equation}
J_n:=P\Biggl(\Biggl|\sum_{i=1}^n\xi_n(X_i,u)-\zeta_n(u)\Biggr|>nu_n^2\omega_1(u_n)\Biggr)\le J_{1n\varepsilon}+J_{2n\varepsilon},
\label{eq:44}
\end{equation}
where
\begin{equation}
J_{1n\varepsilon}=P\Biggl(\Biggl|\sum_{i=1}^n\xi_{n\varepsilon}(X_i,u)-\zeta_n(u)\Biggr|>\tfrac12 nu_n^2\omega_1(u_n)\Biggr)
\label{eq:45}
\end{equation}
and
\begin{equation}
J_{2n\varepsilon}=P\Biggl(\Biggl|\sum_{i=1}^n\bar\xi_{n\varepsilon}(X_i,u)\Biggr|>\tfrac12 nu_n^2\omega_1(u_n)\Biggr).
\label{eq:46}
\end{equation}

Estimate $J_{2n\varepsilon}$. We have
\begin{equation}\label{qqur1}
\begin{split} &
    P\Biggl(\sum_{i=1}^n\bar\xi_{n\varepsilon}(X_i,u)>\tfrac12 nu_n^2\omega_1(u_n)\Biggr)\\
&\le\exp\bigl\{-\gamma_n nu_n^2\omega_1(u_n)/4\bigr\}\bigl(\mathbb{E}\exp\{\gamma_n\bar\xi_{n\varepsilon}(X_1,u)/2\}\bigr)^n\\
&=\exp\bigl\{-\gamma_n nu_n^2\omega_1(u_n)/4\bigr\}\bigl(1+\mathbb{E}\bigl[(\exp\{\gamma_n\xi_n/2\}-1)\mathbf{1}(|\xi_n(X_1,u)|>\varepsilon)\bigr]\bigr)^n\\
&\le\exp\bigl\{-\gamma_n nu_n^2\omega_1(u_n)/4\bigr\}\\
&\qquad\times\bigl(1+\exp\{-\gamma_n\varepsilon/2\}\mathbb{E}\bigl[(\exp\{\gamma_n\xi_n(X_1,u)\}-\exp\{\gamma_n\xi_n(X_1,u)/2\})\mathbf{1}(|\xi_n|>\varepsilon)\bigr]\bigr)^n\\
&\le\exp\bigl\{-\gamma_n nu_n^2\omega_1(u_n)/4+Cn\exp\{-\gamma_n\varepsilon/2\}\bigr\}\\
&=\exp\bigl\{-\gamma_n nu_n^2\omega_1(u_n)/4\,(1+O(1))\bigr\},
\end{split}
\end{equation}
where we used condition B2 in the last equality.

Take $\gamma_{1n}=\gamma_n-2$. Denote
\[
\tilde\xi_{n\varepsilon}(X_1,u)=\log\frac{f(X_1,\theta_0)}{f(X_1,\theta_0+u)}\mathbf{1}\Bigl(\Bigl|\log\frac{f(X_1,\theta_0)}{f(X_1,\theta_0+u)}\Bigr|>\varepsilon\Bigr).
\]
Analogously to the previous estimate one obtains
\begin{equation}\label{qqur2}
\begin{split} 
&P\Biggl(\sum_{i=1}^n\bar\xi_{n\varepsilon}(X_i,u)<-\tfrac12 nu_n^2\omega_1(u_n)\Biggr)
=P\Biggl(\sum_{i=1}^n\tilde\xi_{n\varepsilon}(X_i,u)>\tfrac12 nu_n^2\omega_1(u_n)\Biggr)\\
&\le\exp\bigl\{-\gamma_{1n} nu_n^2\omega_1(u_n)/4\bigr\}\bigl(1+\mathbb{E}\bigl[(\exp\{\gamma_{1n}\tilde\xi_n\}-1)\mathbf{1}(|\tilde\xi_n|>\varepsilon)\bigr]\bigr)^n\\
&\le C\exp\bigl\{-\gamma_n nu_n^2\omega_1(u_n)/4\bigr\}\bigl(1+\mathbb{E}_{\theta_0+u}\bigl[(\exp\{\gamma_n\tilde\xi_n\}-1)\mathbf{1}(|\tilde\xi_n|>\varepsilon)\bigr]\bigr)^n\\
&\le\exp\bigl\{-\gamma_n nu_n^2\omega_1(u_n)/4+Cn\exp\{-\gamma_n\varepsilon/2\}\bigr\}\\
&=\exp\bigl\{-\gamma_n nu_n^2\omega_1(u_n)/4\,(1+o(1))\bigr\}.
\end{split}
\end{equation}

Estimate $J_{1n\varepsilon}$. From Lemma~3.4 in~\cite{3} the estimates
\begin{equation}
\mathbb{E}\bigl[\xi_{n\varepsilon}(X_1,u)-2u^T\phi_\varepsilon(X_1)+\tfrac12 u^T I(\theta_0)u\bigr]=O(u_n^2\omega(u_n))
\label{eq:49}
\end{equation}
and
\begin{equation}
\mathbb{E}\bigl[\xi_{n\varepsilon}(X_1,u)-2u^T\phi_\varepsilon(X_1)+\tfrac12 u^T I(\theta_0)u\bigr]^2=O(u_n^2\omega(u_n))
\label{eq:50}
\end{equation}
follow.

Since
\begin{equation}
|\xi_{n\varepsilon}(X_1,u)-2u^T\phi_\varepsilon(X_1)+\tfrac12 u^T I(\theta_0)u|\le 4\varepsilon
\label{eq:51}
\end{equation}
and relations~\eqref{eq:49},~\eqref{eq:50} hold, applying Prokhorov's inequality~\cite{15} we obtain
\begin{equation}
J_{1n\varepsilon}\le\exp\Biggl\{-\frac{nu_n^2\omega_1(u_n)}{8\varepsilon}\log\Biggl\{\frac{nu_n^2\cdot4\varepsilon}{Cnu_n^2\omega(u_n)}\Biggr\}\Biggr\}
=\exp\Biggl\{\frac{nu_n^2\omega_1(u_n)}{8\varepsilon}\log\frac{\omega(u_n)}{4\varepsilon}\Biggr\}.
\label{eq:52}
\end{equation}
Then we get
\begin{equation}
P\Biggl(\sum_{i=1}^n\xi_{n\varepsilon}(X_i,u)-\zeta_n(u)>nu_n^2\omega_1(u_n)\Biggr)<C\exp\{-nu_n^2/\omega_2(u_n)\},
\label{eq:53}
\end{equation}
provided we choose $\omega_1(u_n)$ so that $\gamma_n\omega_1(u_n)\to\infty$ and $\omega_1(u_n)|\log\omega(u_n)|\to\infty$ as $n\to\infty$, and set
\[
\omega_2(u_n)^{-1}=\min\bigl\{\gamma_n\omega_1(u_n),\ \omega_1(u_n)|\log\omega(u_n)|\bigr\}.
\]
The proof of the inequality
\begin{equation}
P\Biggl(\sum_{i=1}^n\xi_{n\varepsilon}(X_i,u)-\zeta_n(u)<-nu_n^2\omega_1(u_n)\Biggr)<C\exp\{-nu_n^2/\omega_2(u_n)\}
\label{eq:54}
\end{equation}
is analogous and is omitted. It is enough to note that
\begin{equation}
\operatorname{Var}_{\theta_0}\log\frac{f(X_i,\theta_0+u)}{f(X_i,\theta_0)}=\operatorname{Var}_{\theta_0}\log\frac{f(X_i,\theta_0)}{f(X_i,\theta_0+u)}.
\label{eq:55}
\end{equation}
From~\ref{qqur1},~\ref{qqur2} and~\eqref{eq:53},~\eqref{eq:54} relation~\eqref{eq:13} follows.

We now prove~\eqref{eq:13} under condition B1. Take a sequence $\lambda_n$ with $\lambda_n\to\infty$ as $n\to\infty$ such that $\varepsilon_n:=\lambda_n^{-1}|\log u_n|\to0$ and $\gamma_n/\lambda_n\to\infty$ as $n\to\infty$. Choose $\omega_1(u_n)\to0$ so that $\lambda_n\omega_1(u_n)\to\infty$ as $n\to\infty$.

We have
\begin{equation}
J_n\le J_{1n\varepsilon_n}+J_{2n\varepsilon_n}.
\label{eq:56}
\end{equation}
The analogues of the estimates~\ref{qqur1} and~\ref{qqur2} for the upper bound of $J_{2n\varepsilon_n}$ differ only slightly. We give only the analogue of~\ref{qqur1} in the fragment where any differences appear:
\begin{align}
J_{2n\varepsilon_n}
&\le\exp\{-\gamma_n nu_n^2\omega_1(u_n)/4\}\bigl(1+\exp\{-\gamma_n\varepsilon_n/2\}\notag\\
&\qquad\times\mathbb{E}\bigl[(\exp\{\gamma_n\xi_n(X_i,u)\}-\exp\{\gamma_n\xi_n(X_i,u)/2\})\mathbf{1}(|\xi_n(X_i,u)|>\varepsilon_n)\bigr]\bigr)^n\notag\\
&\le\exp\{-\gamma_n nu_n^2\omega_1(u_n)/4+Cn\exp\{-\gamma_n\varepsilon_n/2\}\}\notag\\
&=\exp\{-\gamma_n nu_n^2\omega_1(u_n)/4\,(1+O(1))\}.
\label{eq:57}
\end{align}
In order to estimate $J_{1n\varepsilon_n}$ by means of Prokhorov's inequality we need to show that
\begin{equation}
\mathbb{E}\bigl[\xi_{n\varepsilon_n}(X_1,u)-2u^T\phi_\varepsilon(X_1)+\tfrac12 u^T I(\theta_0)u\bigr]=O(u_n^2\omega(u_n))
\label{eq:58}
\end{equation}
and
\begin{equation}
\mathbb{E}\bigl[\xi_{n\varepsilon_n}(X_1,u)-2u^T\phi_\varepsilon(X_1)+\tfrac12 u^T I(\theta_0)u\bigr]^2=O(u_n^2\omega(u_n)).
\label{eq:59}
\end{equation}
To prove~\eqref{eq:58} and~\eqref{eq:59} it is sufficient to show that
\begin{equation}
\mathbb{E}\bigl[\bar\xi_{n\varepsilon}(X_1,u)\mathbf{1}(|\xi_n(X_1,u)|<\varepsilon_n)\bigr]=O(u_n^2\omega(u_n))
\label{eq:60}
\end{equation}
and
\begin{equation}
\mathbb{E}\bigl[\bar\xi_{n\varepsilon}^2(X_1,u)\mathbf{1}(|\xi_n(X_1,u)|<\varepsilon_n)\bigr]=O(u_n^2\omega(u_n)).
\label{eq:61}
\end{equation}
We have
\begin{equation}
\begin{split} &
\mathbb{E}\bigl[\bar\xi_{n\varepsilon}(X_1,u)\mathbf{1}(\xi_n(X_1,u)>\varepsilon)\bigr]\le\mathbb{E}\bigl[g_n(X_1,u)\mathbf{1}(g_n(X_1,u_n)>\varepsilon)\bigr] \\ &
\le\varepsilon^{-1}\mathbb{E}\bigl[g_n^2(X_1,u)\mathbf{1}(g_n(X_1,u_n)>\varepsilon)\bigr]=O(u_n^2\omega(u_n))
\label{eq:62}
\end{split}
\end{equation}
and
\begin{equation}
\mathbb{E}\bigl[\bar\xi_{n\varepsilon}^2(X_1,u)\mathbf{1}(\xi_n(X_1,u)>\varepsilon)\bigr]
\le\mathbb{E}\bigl[g_n^2(X_1,u)\mathbf{1}(g_n(X_1,u_n)>\varepsilon)\bigr]=O(u_n^2\omega(u_n)),
\label{eq:63}
\end{equation}
where the last equalities in~\eqref{eq:62} and~\eqref{eq:63} were proved in~(3.5) of~\cite{3}.

From~\eqref{eq:62},~\eqref{eq:63} and condition C1 relations~\eqref{eq:60} and~\eqref{eq:61} follow, and therefore also~\eqref{eq:58} and~\eqref{eq:59}. Using~\eqref{eq:58},~\eqref{eq:59} and applying Prokhorov's inequality we obtain
\begin{equation}
J_{1n\varepsilon_n}\le\exp\Biggl\{-\frac{nu_n^2\omega_1(u_n)}{8\varepsilon_n}\log\Biggl\{\frac{nu_n^2\cdot4\varepsilon_n}{Cnu_n^{2+\gamma}}\Biggr\}\Biggr\}
\le\exp\{-Cnu_n^2\omega_1(u_n)\lambda_n(1+o(1))\}.
\label{eq:64}
\end{equation}
From~\eqref{eq:57} and~\eqref{eq:64} relation~\eqref{eq:13} follows.

\subsection{Estimates of the rate of uniform approximation}
\label{sub:uniform}

Denote
\begin{equation}
Z_{n\theta}(u)=\prod_{j=1}^n\frac{f(X_j,\theta+u)}{f(X_j,\theta)}.
\label{eq:65}
\end{equation}
We obtain estimates of the moduli of continuity of $Z_n^{1/2}(u)$ and of the log-likelihood ratio.

We begin with the proof of~\eqref{eq:14}. On the set $[-Cu_n,Cu_n]^d$ take a $\delta_n$-net consisting of the points $u_{ni}$, $1\le i\le l_n$, where the value of $\delta_n$ will be defined later.

For the proof of~\eqref{eq:14} we consider two cases:
\begin{equation}
u_n>cn^{-1/2}\sqrt{\log n},\qquad c>1
\label{eq:66}
\end{equation}
and
\begin{equation}
u_n\le cn^{-1/2}\sqrt{\log n},\qquad c<1.
\label{eq:67}
\end{equation}

We start with the case~\eqref{eq:66}. Take a sequence $\omega_3(u_n)$ such that $\omega_3(u_n)\to0$ and $\omega_2(u_n)/\omega_3(u_n)\to\infty$ as $n\to\infty$. Put
\[
\delta_n=\exp\{-nu_n^2/\omega_3(u_n)\}.
\]
Define the event
\begin{equation}
A_n=\Biggl\{X_1,\dots,X_n:\max_{1\le i\le l_n}\Biggl|\sum_{i=1}^n\xi_n(X_i,u_{ni})-\zeta_n(u_{ni})\Biggr|<nu_n^2\omega_1(u_n)/2\Biggr\}.
\label{eq:68}
\end{equation}
We have
\begin{align}
P(\overline{A}_n)
&\le\sum_{i=1}^{l_n}P\Biggl(\Biggl|\sum_{i=1}^n\xi_n(X_i,u_{ni})-\zeta_n(u_{ni})\Biggr|>nu_n^2\omega_1(u_n)/2\Biggr)\notag\\
&\le l_n\exp\{-nu_n^2/\omega_2(u_n)\}\le\exp\{-nu_n^2(1/\omega_2(u_n)+C/\omega_3(u_n))\}\notag\\
&=\exp\{-nu_n^2/\omega_2(u_n)\,(1+o(1))\}.
\label{eq:69}
\end{align}
For the proof of~\eqref{eq:14} we use the following lemma (Lemma~3.1 of Chapter~3 in~\cite{10}).

\begin{lemma}\label{lem:4.2}
Let $\xi_1,\dots,\xi_n$ be independent identically distributed random variables with $\mathbb{E}\xi_1=0$ and $\mathbb{E}|\xi_1|^{2\beta}<\infty$, $\beta\ge1$. Then
\begin{equation}
\mathbb{E}\Biggl|\sum_{i=1}^n\xi_i\Biggr|^{2\beta}\le C_\beta n^\beta\mathbb{E}|\xi_1|^{2\beta}.
\label{eq:70}
\end{equation}
\end{lemma}

Hence, using condition E, we obtain
\begin{align}
&\mathbb{E}\Biggl|\sum_{i=1}^n\bigl(\log f(X_i,\theta_0+u)-\log f(X_i,\theta+v)-(u-v)^T\phi_\varepsilon(X_i)\bigr)\Biggr|^{\beta_1}\notag\\
&\le Cn^{\beta_1/2}\mathbb{E}\bigl(\xi_n(X_1,u,v-u)-(u-v)^T\phi_\varepsilon(X_1)\bigr)^{\beta_1}\le Cn^{\beta_1/2}|u-v|^{\beta_2}.
\label{eq:71}
\end{align}
Applying Theorem~19 of Appendix~1 in~\cite{10} we then obtain
\begin{align}
&P\Biggl(\sup_i\sup_{|u-u_{in}|<\delta_n}\Biggl|\sum_{s=1}^n\xi_n(X_s,u)-\zeta_n(u)\Biggr|>nu_n^2\omega_3(u_n)\Biggr)\notag\\
&\le(nu_n^2\omega_3(u_n))^{-1}\mathbb{E}\Biggl(\sup_i\Biggl|\sum_{s=1}^n\xi_n(X_s,u_{in})-\zeta_n(u_{in})\Biggr|\Biggr)\notag\\
&\quad+(nu_n^2\omega_3(u_n))^{-1}\mathbb{E}\Bigl[\sup_i\sup_{|u-u_{in}|<\delta_n}|\zeta_n(u)-\zeta_n(u_{in})|\Bigr]\notag\\
&\quad+(u_n^2\omega_3(u_n))^{-1}\delta_n\mathbb{E}|\phi_\varepsilon(X_1)|\notag\\
&\quad+(nu_n^2\omega_3(u_n))^{-1}\mathbb{E}\Biggl(\sup_i\sup_{|u-u_{in}|<\delta_n}\Biggl|\sum_{s=1}^n\bigl(\xi_n(X_s,u_{in},u-u_{in})+(u_{in}-u)\phi(X_s)\bigr)\Biggr|\Biggr)\notag\\
&\le C(nu_n^2\omega_3(u_n))^{-1}(\delta_n+\delta_n^{(\beta_2-d)/\beta}).
\label{eq:72}
\end{align}
From~\eqref{eq:69} and~\eqref{eq:72} we obtain~\eqref{eq:14} in the case~\eqref{eq:66}. In the case~\eqref{eq:67} the estimates for the proof of~\eqref{eq:14} are analogous and are omitted.

We now estimate the modulus of continuity of $Z_n^{1/2}(u)$ when $\Theta\subset\mathbb{R}$. Take a sequence $\omega_3(u_n)$ such that $\omega_3(u_n)\to0$ and $\omega_2(u_n)/\omega_3(u_n)\to\infty$ as $n\to\infty$. Put
\[
\delta_n=\exp\{-nu_n^2/\omega_3(u_n)\}.
\]
Using the estimate of Lemma~5.5 of Chapter~1 in~\cite{10} we obtain
\begin{equation}
\mathbb{E}_{\theta_0}\bigl(Z_n^{1/2}(u)-Z_n^{1/2}(v)\bigr)^2\le Cn|u-v|^2
\label{eq:73}
\end{equation}
for $|u|<C$, $|v|<C$ and $n>n_0$.

Take a sequence $\omega_4(u_n)$ such that $\omega_4(u_n)\to0$ and $\omega_4(u_n)/\omega_3(u_n)\to\infty$ as $n\to\infty$. In the case~\eqref{eq:66}, from~\eqref{eq:73} and Theorem~19 of the Appendix in~\cite{10} we obtain
\begin{align}
&P\Biggl(\sup_i\sup_{|u-u_{in}|<\delta_n}\bigl|Z_n^{1/2}(u)-Z_n^{1/2}(u_{in})\bigr|>\exp\{-nu_n^2/\omega_4(u_n)\}\Biggr)\notag\\
&\le\exp\{nu_n^2/\omega_4(u_n)\}\mathbb{E}\Bigl[\sup_i\sup_{|u-u_{in}|<\delta_n}\bigl|Z_n^{1/2}(u)-Z_n^{1/2}(u_{in})\bigr|\Bigr]\notag\\
&\le Cn^{1/2}\exp\{nu_n^2/\omega_4(u_n)\}\delta_n^{1/2}\le C\exp\Bigl\{-\frac{nu_n^2}{2\omega_3(u_n)}\Bigr\}.
\label{eq:74}
\end{align}

In the multidimensional case, for the proof of the analogue of inequality~\eqref{eq:74} one uses, instead of~\eqref{eq:73}, the inequality
\begin{equation}
\mathbb{E}_\theta\bigl(Z_{n\theta}^{1/m}(u)-Z_{n\theta}^{1/m}(v)\bigr)^m\le Cn^{m/2}|u-v|^m
\label{eq:75}
\end{equation}
(see Lemma~3.2 of Chapter~2 in~\cite{10}). It holds provided condition D is satisfied. The remaining estimates are analogous.

We obtain an estimate of the modulus of continuity of $Z_n^{1/2}(u)$ in the case~\eqref{eq:67} when $\Theta\subset\mathbb{R}$. Put $\delta_n=\alpha_n u_n$, where $\alpha_n=n^{-c_1}$ with $c_1>1/2+8c^2$. Then
\begin{equation}
P(\overline{A}_n)\le\exp\{-nu_n^2/\omega_2(u_n)-\log\alpha_n\}=\exp\{-nu_n^2/\omega_2(u_n)\,(1+o(1))\}.
\label{eq:76}
\end{equation}
Arguing as in~\eqref{eq:72} we obtain
\begin{equation}
\begin{split} &
P\Biggl(\sup_i\sup_{|u-u_{in}|<\delta_n}\bigl|Z_{n\theta}^{1/2}(u)-Z_{n\theta}^{1/2}(u_{in})\bigr|>\exp\{-2nu_n^2\}\Biggr)\\ &
\le Cn^{1/2}\alpha_n^{1/2}u_n^{1/2}\exp\{2nu_n^2\}=O\bigl(n^{1/4-c_1/2+4c^2}\log^{1/4}n\bigr).
\label{eq:77}
\end{split}
\end{equation}
Note that the estimates~\eqref{eq:76} and~\eqref{eq:77} are valid for any $c>1$ and $c_1>1/2+8c^2$. This allows one to prove their validity in the multidimensional case for $c>m/2$ and $c_1>m(m/2+d/4+2c^2)/(m-d)$, using condition D.

\subsection{Proof of Theorem~\ref{th:MLE}}
\label{sub:MLE}

Thus we have shown that the log-likelihood ratio, regarded as a random process, is well approximated by the random process $\zeta_n(u)$, $|u|<Cu_n$. Consequently the proof of the theorem is essentially reduced to the study of the point at which the maximum of the process $\zeta_n(u)$, $|u|<Cu_n$, is attained.

We say that an estimator $\tilde\theta_n$ is $u_n$-consistent if for every $\theta\in\Theta$ there exists a neighbourhood $U\subset\Theta$ of $\theta$ such that for every $\delta>0$
\[
\lim_{n\to\infty}\sup_{\theta\in U}P_\theta(|\tilde\theta_n-\theta|>\delta u_n)=0.
\]
If an estimator is $u_n$-consistent and the Fisher information is finite, then the lower bound of asymptotic efficiency in the Bahadur sense holds for it in the zone of moderate deviation probabilities (see Theorem~2.2 of~\cite{5}). This lower bound coincides with the lower bound of the moderate-deviation probabilities of the maximum-likelihood estimator in Theorem~\ref{th:MLE}. If the conditions of Theorem~3.2 of Chapter~2 in~\cite{10} are satisfied, then the maximum-likelihood estimator is asymptotically normal and therefore $u_n$-consistent. If the conditions of Theorem~\ref{th:MLE} are satisfied, then the conditions of Theorem~3.2 of Chapter~2 in~\cite{10} are also satisfied. Thus, to prove~\eqref{eq:18} it remains only to prove the upper bound.

Define the event
\begin{align*}
D_n=\bigl\{X_1,\dots,X_n:&\ |Z_n(u)-Z_n(u_{in})|\le C\exp\{-2nu_n^2\},\\
&\ \exp\{\zeta_n(u_{in})-nu_{in}^2\omega_1(u_n)\}<Z_n(u_{in})<\exp\{\zeta_n(u_{in})+u_{in}^2\omega_1(u_n)\},\\
&\ |u-u_{in}|<\delta_n,\ 1\le i\le l_n\bigr\}.
\end{align*}
From~\eqref{eq:69},~\eqref{eq:76} and~\eqref{eq:74},~\eqref{eq:77} it follows that
\begin{equation}
P(\overline{D}_n)\le C\exp\{-C_1 nu_n^2\},
\label{eq:78}
\end{equation}
where the constant $C_1$ can be taken arbitrarily large for a suitable $C=C(C_1)$.

To prove the upper bound we show that if $0<r<1$, then for the event
\[
\Upsilon_n(r)=\{X_1,\dots,X_n:\psi_n^T\psi_n<nr u_n^2\}
\]
one has
\begin{equation}
P\bigl(|I^{1/2}(\theta)(\hat\theta_n-\theta)|>u_n,\ D_n\cap\Upsilon_n(r)\bigr)=0.
\label{eq:79}
\end{equation}
By Theorem~\ref{th:LD}, for the complementary events $\overline{\Upsilon}_n(r)$ one has
\begin{equation}
\lim_{n\to\infty}(nu_n^2)^{-1}\log P(\overline{\Upsilon}_n(r))=-r/2.
\label{eq:80}
\end{equation}
Consequently from~\eqref{eq:78}--\eqref{eq:80} the upper bound in~\eqref{eq:18} follows.

Note that if the event $D_n$ holds, then the events $|\psi_n|<n^{1/2}u_n(1-\varepsilon)/2$ ($0<\varepsilon<1$) and $|I^{1/2}(\theta_0)(\hat\theta_n-\theta_0)|>u_n$ cannot hold simultaneously. They would simultaneously imply that $Z_n(\hat\theta_n)<1$ and $Z_n(0)=1$, which is a contradiction.

Therefore in what follows we shall assume that the event
\begin{equation}
G_n=\{X_1,\dots,X_n:|\psi_n|>n^{1/2}u_n/2\}
\label{eq:81}
\end{equation}
holds. Take $\kappa$ with $0<\kappa<r$ and define the events
\begin{equation}
W_n(\kappa)=D_n\cap\{X_1,\dots,X_n:|\psi_n|<(r-\kappa)n^{1/2}u_n\}.
\label{eq:82}
\end{equation}
Observe that
\begin{align}
\max_{u^TI(\theta_0)u<(r-\kappa)u_n^2}\zeta_n(u)
&=-\tfrac12\min_{u^TI(\theta_0)u<(r-\kappa)u_n^2}|n^{-1/2}\psi_n-\sqrt{n}\,u^TI^{1/2}(\theta_0)|^2\notag\\
&\qquad+\tfrac12 n\psi_n^T\psi_n=\tfrac12 n\psi_n^T\psi_n.
\label{eq:83}
\end{align}
Assume that the case~\eqref{eq:66} holds. Then on the set of events $W_n(\kappa)$, for a point $u_{in}$ with $|u_{in}-\hat\theta_n|<\delta_n$, one has
\begin{align}
\bigl|\log Z_n^{1/2}(\hat\theta_n)-\log Z_n^{1/2}(u_{in})\bigr|
&\le\bigl|\log\bigl(1+Z_n^{-1/2}(u_{in})(Z_n^{1/2}(\hat\theta_n)-Z_n^{1/2}(u_{in}))\bigr)\bigr|\notag\\
&\le C Z_n^{-1/2}(u_{in})\bigl(Z_n^{1/2}(\hat\theta_n)-Z_n^{1/2}(u_{in})\bigr)\notag\\
&\le C\exp\{-cnu_n^2/\omega_3(u_n)\},
\label{eq:84}
\end{align}
where the second inequality follows from the fact that $Z_n^{-1/2}(u_{in})(Z_n^{1/2}(\hat\theta_n)-Z_n^{1/2}(u_{in}))$ is positive (since $\hat\theta_n$ is a maximum-likelihood estimator) and the third inequality follows from the occurrence of the event $D_n$.

Consider the case~\eqref{eq:67}. Estimating analogously to~\eqref{eq:84} and applying the same arguments we obtain
\begin{align}
\bigl|\log Z_n^{1/2}(\hat\theta_n)-\log Z_n^{1/2}(u_{in})\bigr|
&\le C Z_n^{-1/2}(u_{in})\bigl(Z_n^{1/2}(\hat\theta_n)-Z_n^{1/2}(u_{in})\bigr)\notag\\
&\le C\exp\bigl\{-cnu_n^2/\omega_3(u_n)\,(2-r)/2\bigr\}.
\label{eq:85}
\end{align}
From~\eqref{eq:84} and~\eqref{eq:85} it follows that on the set of events $D_n$, for a point $u_{in}$ with $|u_{in}-\hat\theta_n|<\delta_n$, one has
\begin{equation}
\bigl|\log Z_n^{1/2}(\hat\theta_n)-\log Z_n^{1/2}(u_{in})\bigr|\le o(nu_n^2).
\label{eq:86}
\end{equation}
Therefore on the set of events $D_n\cap W_n(\kappa)$ one has
\begin{equation}
\max_{u^TI(\theta_0)u<(r-\kappa)u_n^2}\log Z_n(u)=\tfrac12 n\psi_n^T\psi_n\,(1+o(1)),
\label{eq:87}
\end{equation}
and the maximum is attained in the region $\{u:u^TI(\theta_0)u<r^2 u_n^2\}$, which completes the proof of~\eqref{eq:18}.

Let $e\in\mathbb{R}^d$ with $|e|=1$. Denote by $B_\delta$ the ball of radius $\delta$, $0<\delta<0.05$, centred at the origin.

To prove~\eqref{eq:21}--\eqref{eq:23} it is sufficient to show that
\begin{equation}
\lim_{n\to\infty}(nu_n^2)^{-1}\log P_\theta\bigl(I^{1/2}(\theta_0)\psi_n\in u_nn^{1/2}(e+B_\delta),\ \hat\theta_n-\theta_0\notin u_n(e+B_{4\sqrt{\delta}})\bigr)=-\infty.
\label{eq:88}
\end{equation}
For notational simplicity we carry out the argument for the proof of~\eqref{eq:88} when the matrix $I(\theta_0)$ is the identity. It is enough to prove~\eqref{eq:88} under the event $D_n$. Then, if $\psi_n\in n^{1/2}u_n(e+B_\delta)$, the smallest possible value of $\max\{\zeta_n(u):u\in\mathbb{R}^d\}$ equals $(1-\delta)^2 nu_n^2/2$. At the same time the largest possible value of $\max\{\zeta_n(u):u\notin u_n(e+B_{4\sqrt{\delta}})\}$ equals $nu_n^2\bigl((1+\delta)^2-8\delta\bigr)/2$. Consequently~\eqref{eq:88} holds.

\section{Proof of Theorem~\ref{th:Bayes}}

Since under the conditions of Theorem~\ref{th:Bayes} the Bayesian estimator is asymptotically normal (see Theorem~3.1 of Chapter~3 in~\cite{10}), it is $u_n$-consistent. Consequently the lower bound in~\eqref{eq:33} follows from the lower bound for asymptotic efficiency in the Bahadur sense in the zone of moderate deviation probabilities (see Theorem~2.2 of~\cite{5}).

The proof of the upper bound in~\eqref{eq:33} will be based on the same approach as the proof of Theorem~\ref{th:MLE}. We define an event $D_n$ such that $(nu_n^2)^{-1}\log P_{\theta_0}(\overline{D}_n)\to-\infty$ as $n\to\infty$, and show that for every $\delta>0$
\begin{equation}
P_{\theta_0}\bigl(|I^{1/2}(\theta_0)(\tilde\theta_n-\theta_0)|>u_n,\ |\psi_n|<n^{1/2}(1-\delta)u_n,\ D_n\bigr)=0.
\label{eq:89}
\end{equation}
Then the upper bound will follow from Theorem~\ref{th:LD}.

For $r>0$ denote
\begin{equation}
I(r)=\int_{|u|>ru_n}Z_n(u)\pi(u)\,du,
\label{eq:90}
\end{equation}
\begin{equation}
J(r)=\int_{|u|>ru_n}|u|^a Z_n(u)\pi(u)\,du,\qquad a>0,
\label{eq:91}
\end{equation}
and
\begin{equation}
Q(r)=I(r)\Biggl(\int Z_n(u)\pi(u)\,du\Biggr)^{-1}.
\label{eq:92}
\end{equation}
From Lemma~5.2 of Chapter~1 and Theorem~5.5 of Chapter~1 in~\cite{10} it follows that for any $a$ one can choose sequences $\omega_1(u_n)\to0$, $\omega_2(u_n)\to0$ (satisfying~\eqref{eq:13}) and a sequence $r_n$ with $r_nu_n\to0$, $r_n\to\infty$ as $n\to\infty$ such that for $n>n_0$
\begin{align}
P\bigl(I(r_n)>\exp\{-nu_n^2/\omega_1(u_n)\}\bigr)&\le\exp\{-nu_n^2/\omega_2(u_n)\},
\label{eq:93}\\
P\bigl(J(r_n)>\exp\{-nu_n^2/\omega_1(u_n)\}\bigr)&\le\exp\{-nu_n^2/\omega_2(u_n)\},
\label{eq:94}\\
P\bigl(Q(r_n)>\exp\{-nu_n^2/\omega_1(u_n)\}\bigr)&\le\exp\{-nu_n^2/\omega_2(u_n)\}.
\label{eq:95}
\end{align}
Moreover there exists $\omega_3(u_n)\to0$ as $n\to\infty$ such that
\begin{equation}
P(\tilde\theta_n-\theta_0>r_nu_n/2)\le\exp\{-nu_n^2/\omega_3(u_n)\}.
\label{eq:96}
\end{equation}
This allows us, in the subsequent estimates, to assume that the complementary events to those appearing on the left-hand sides of~\eqref{eq:93}--\eqref{eq:96} hold, and to carry out the further arguments for the Bayesian estimator $\tilde\theta_n$ defined not by equation~\eqref{eq:Bayes} but by
\begin{equation}
\tilde\theta_n=\arg\min_{t\in\Theta}\frac{\displaystyle\int_{|\theta-\theta_0|<r_nu_n}l(t-\theta)\prod_{j=1}^n f(X_j,\theta)\pi(\theta)\,d\theta}{\displaystyle\int_{|\theta-\theta_0|<r_nu_n}\prod_{j=1}^n f(X_j,\theta)\pi(\theta)\,d\theta},
\label{eq:97}
\end{equation}
and, in addition, to assume that the minimizer $\tilde\theta_n$ lies in the region $|\tilde\theta-\theta_0|<r_nu_n/2$.

We carry out the subsequent arguments for the case
\begin{equation}
u_n\sqrt{\frac n{\log n}}\to\infty\qquad(n\to\infty)
\label{eq:98}
\end{equation}
and afterwards indicate the differences for the case
\begin{equation}
u_n\sqrt{\frac n{\log n}}<c,
\label{eq:99}
\end{equation}
where $c$ is an arbitrary positive constant.

Instead of estimates of the modulus of continuity and the application of Theorem~19 of Appendix~1 in~\cite{10}, our estimates will be based on the following auxiliary inequality:
\begin{align}
\mathbb{E}|Z_n(u)-Z_n(v)|
&=\mathbb{E}\bigl|(Z_n^{1/2}(u)-Z_n^{1/2}(v))(Z_n^{1/2}(u)+Z_n^{1/2}(v))\bigr|\notag\\
&\le\bigl(\mathbb{E}(Z_n^{1/2}(u)-Z_n^{1/2}(v))^2\bigr)^{1/2}\bigl(\mathbb{E}(Z_n^{1/2}(u)+Z_n^{1/2}(v))^2\bigr)^{1/2}\notag\\
&\le2\bigl(\mathbb{E}(Z_n^{1/2}(u)-Z_n^{1/2}(v))^2\bigr)^{1/2}\le Cn|u-v|.
\label{eq:100}
\end{align}
Here in the last inequality we used Lemma~5.5 of Chapter~1 in~\cite{10}.

Take a sequence $r_n>0$ with $r_n\to\infty$, $r_nu_n\to0$ as $n\to\infty$. On the set $K_n=(-r_nu_n,r_nu_n)^d$ construct a $\delta_n$-net consisting of the points $u_{ni}$, $1\le i\le l_n$. The choice of the sequence $\delta_n>0$ will be different for the cases~\eqref{eq:98} and~\eqref{eq:99}. Partition the set $K_n$ into measurable pairwise disjoint subsets $\Delta_{ni}$, $u_{ni}\in\Delta_{ni}$, $1\le i\le k_n$, of diameter $\sup_{u,v\in\Delta_{ni}}|u-v|<3\delta_n$.

Define the event $A_n$ by~\eqref{eq:68}. Denote
\begin{equation}
L_n(t)=\int_{K_n}Z_n(u)\,du-\sum_{i=1}^{k_n}\int_{\Delta_{ni}}Z_n(u_{ni})\,du.
\label{eq:101}
\end{equation}
Using~\eqref{eq:100} we obtain
\begin{equation}
J_n=\mathbb{E}|L_n|\le\sum_{i=1}^{k_n}\int_{\Delta_{ni}}\mathbb{E}|Z_n(u)-Z_n(u_{ni})|\,du\le Cn\delta_n(r_nu_n)^d.
\label{eq:102}
\end{equation}
Note that in~\eqref{eq:102} the integration over the cube $K_n$ may be replaced by integration over any union of the sets $\Delta_{ni}$.

Without loss of generality we assume in the arguments that
\[
\inf\bigl\{\|u\|:u^TI(\theta_0)u=1,\ u\in\mathbb{R}^d\bigr\}=1.
\]
Denote $C_0=\sup\bigl\{\|u\|:|u^TI(\theta_0)u|=1,\ u\in\mathbb{R}^d\bigr\}$.

We begin with the case~\eqref{eq:98}. Let $\|t\|>u_n$. For $\|u\|<u_n/2$ we have $\|t-u\|\ge\|t\|-\|u\|>\|u\|$. Consequently, for $0<r_1<1/2$,
\begin{equation}
\int_{\|u\|<r_1u_n}l(t-u)Z_n(u)\,du\ge\int_{\|u\|<r_1u_n}l(u)Z_n(u)\,du.
\label{eq:103}
\end{equation}
For $0<\delta_0<0.1C_0^{-1}$ consider the case when the event
\begin{equation}
\Upsilon_n=\{X_1,\dots,X_n:|\psi_n|<\delta_0 n^{1/2}u_n\}
\label{eq:104}
\end{equation}
holds. Take $r_1,r_2$ such that $0<r_1<C_0r_1<r_2<C_0r_2<1/2$. Then
\[
U_n=\{u:C_0r_1u_n<|uI^{1/2}(\theta_0)|<r_2u_n\}\subset V_n=\{u:r_1u_n\le\|u\|\le r_2u_n\}.
\]
Consequently, by~\eqref{eq:30},
\begin{equation}
\int_{V_n}\bigl(l(t-u)-l(u)\bigr)Z_n(u)\,du\ge C l_1(u_n)\int_{U_n}Z_n(u)\,du.
\label{eq:105}
\end{equation}
Take a sequence $\omega_3(u_n)$ such that $\omega_3(u_n)\to0$ and $\omega_3(u_n)/\omega_2(u_n)\to\infty$ as $n\to\infty$. Put $\delta_n=\exp\{-nu_n^2/\omega_3(u_n)\}$. Denote by $M_n$ the set of points $u_{ni}$, $1\le i\le l_n$, such that $u_{ni}\in U_n$. Denote $W_{1n}=\bigcup_{u_{ni}\in M_n}\Delta_{ni}$.

Put
\begin{equation}
L_{1n}=L_n(r_1,r_2)=\int_{W_{1n}}Z_n(u)\,du-\sum_{u_{ni}\in M_n}\int_{\Delta_{ni}}l(t-u)Z_n(u_{ni})\,du.
\label{eq:106}
\end{equation}
Estimating analogously to~\eqref{eq:102} we obtain
\begin{equation}
\mathbb{E}|L_{1n}|\le Cn\delta_nu_n^d.
\label{eq:107}
\end{equation}
Define the event
\[
D_{1n}=D_n(r_1,r_2)=\bigl\{X_1,\dots,X_n:|L_{1n}|<\exp\{-nu_n^2/(3\omega_3(u_n))\}\bigr\}.
\]
Then, using~\eqref{eq:107},
\begin{equation}
P(\overline{D}_{1n})<\exp\{nu_n^2/(3\omega_3(u_n))\}\mathbb{E}|L_{1n}|\le C\exp\{-2nu_n^2/(3\omega_3(u_n))\}.
\label{eq:108}
\end{equation}
We continue the argument under the assumption that the event $A_n\cap D_{1n}$ holds.

Since the event $A_n$ holds,
\begin{equation}
Z_n(u_{ni})\exp\{-\zeta(u_{ni})\}\ge\exp\{-nu_{ni}^2\omega_1(u_{ni})\}.
\label{eq:109}
\end{equation}
Consequently
\begin{equation}
Z_n(u_{ni})\ge\exp\{-2n|u_{ni}|\delta_0 u_n-nu_{ni}^TI(\theta_0)u_{ni}/2-nu_n^2\omega_1(u_n)\}.
\label{eq:110}
\end{equation}
Hence
\begin{equation}
\sum_{u_{ni}\in M_n}\int_{\Delta_{ni}}Z_n(u_{ni})\,du\ge\exp\bigl\{-n(r_2^2+4\delta_0)u_n^2/2-nu_n^2\omega_1(u_n)(1+o(1))\bigr\}.
\label{eq:111}
\end{equation}
Therefore
\begin{align}
\int_{V_{1n}}\bigl(l(t-u)-l(u)\bigr)Z_n(u)\,du
&\ge l_1(u_n)n^{-d/2}\exp\bigl\{-n(r_2^2+4\delta_0)u_n^2/2-nu_n^2\omega_1(u_n)(1+o(1))\bigr\}.
\label{eq:112}
\end{align}
We have
\begin{equation}
\int_{r_2u_n\le\|u\|<u_n/2}\bigl(l(t-u)-l(u)\bigr)Z_n(u)\,du>0.
\label{eq:113}
\end{equation}
We estimate from above the integral
\begin{equation}
\int_{u_n/2\le\|u\|<r_nu_n}\bigl(l(t-u)-l(u)\bigr)Z_n(u)\,du.
\label{eq:114}
\end{equation}
Denote by $M_{2n}$ a $\delta_n$-net consisting of the points $u_{ni}$, $1\le i\le l_{1n}$, such that $0.5u_n\le\|u_{ni}\|\le r_nu_n$. Define the event $D_{2n}=D_n(0.5,r_n)$.

Analogously to~\eqref{eq:108} we have
\begin{equation}
P(\overline{D}_{2n})\le C\exp\{-2nu_n^2/(3\omega_1(u_n))(1+o(1))\}.
\label{eq:115}
\end{equation}
If the event $A_n$ holds, then for the points $u_{ni}$, $1\le i\le l_{1n}$,
\begin{equation}
Z_n(u_{ni})\exp\{-\zeta_n(u_{ni})\}\le\exp\{n|u_{ni}|^2\omega_1(|u_{ni}|)\}.
\label{eq:116}
\end{equation}
Consequently
\begin{equation}
Z_n(u_{ni})\le\exp\{2n|u_{ni}|\delta_0 u_n-nu_{ni}^TI(\theta_0)u_{ni}/2+nu_n^2\omega(u_n)\}.
\label{eq:117}
\end{equation}
Hence
\begin{equation}
\sum_{u_{ni}\in M_{2n}}\int_{\Delta_{ni}}Z_n(u_{ni})\,du\le\exp\bigl\{-n(1/4-4\delta_0)u_n^2/2+nu_n^2\omega(u_n)(1+o(1))\bigr\}.
\label{eq:118}
\end{equation}
Therefore, if the event $A_n\cap D_{2n}$ holds,
\begin{align}
&\int_{u_n/2\le\|u\|<r_nu_n}\bigl(l(t-u)-l(u)\bigr)Z_n(u)\,du\notag\\
&\le C l_1(r_nu_n)\int_{u_n/2\le\|u\|<r_nu_n}Z_n(u)\,du\notag\\
&\le C l_1(r_nu_n)n^{-d/2}\exp\bigl\{-n(1/8-2\delta_0)u_n^2+nu_n^2\omega_1(u_n)(1+o(1))\bigr\}.
\label{eq:119}
\end{align}
Since the right-hand side of~\eqref{eq:112} is larger than the right-hand side of~\eqref{eq:119}, from~\eqref{eq:103},~\eqref{eq:112},~\eqref{eq:113} and~\eqref{eq:119} relation~\eqref{eq:89} follows, provided the event $\Upsilon_n\cap A_n\cap D_{1n}\cap D_{2n}$ holds.

Take $\delta>0$ with $\delta<\delta_0/10$. Consider the case $\delta_0 n^{1/2}u_n<|\psi_n|<(1-5\delta C_0)n^{1/2}u_n$.

Denote by $B(v,C)=\{x:|I^{1/2}(\theta_0)(x-v)|<C,\ x\in\mathbb{R}^d\}$ the ball centred at $v\in\mathbb{R}^d$ of radius $C$.

Let the points $x_{n1},\dots,x_{nk}$ form a $\delta u_n$-net in the set $B(0,(1-5\delta)u_n)\setminus B(0,\delta_0 u_n)$. Denote by $J_{ni}$ the set of all indices $j$, $1\le j\le l_n$, such that $\Delta_{nj}\subset B(x_{ni},2\delta u_n)$. Denote $\overline{J}_{ni}=\{j:j\notin J_{ni},\ 1\le j\le k_n\}$.

Put $U_{ni}=\bigcap_{j\in J_{ni}}\Delta_{nj}$. Denote $V_{ni}=K_n\setminus U_{ni}$.

For $1\le i\le k$ define
\begin{align}
L_{ni}&=\int_{U_{ni}}Z_n(u)\,du-\sum_{j\in J_{ni}}\int_{\Delta_{nj}}Z_n(u_{ni})\,du,
\label{eq:120}\\
\overline{L}_{ni}&=\int_{V_{ni}}Z_n(u)\,du-\sum_{j\in\overline{J}_{ni}}\int_{\Delta_{nj}}Z_n(u_{ni})\,du.
\label{eq:121}
\end{align}
Define the events
\begin{align}
D_{ni}&=\bigl\{X_1,\dots,X_n:|L_{ni}|<\exp\{-nu_n^2/(3\omega_3(u_n))\}\bigr\},
\label{eq:122}\\
G_{ni}&=\bigl\{X_1,\dots,X_n:|\overline{L}_{ni}|<\exp\{-nu_n^2/(3\omega_3(u_n))\}\bigr\}.
\label{eq:123}
\end{align}
Estimating analogously to~\eqref{eq:102} and~\eqref{eq:108} we obtain
\begin{equation}
P(\overline{D}_{ni})<C\exp\{-2nu_n^2/(3\omega_3(u_n))\},\qquad
P(\overline{G}_{ni})<C\exp\{-2nu_n^2/(3\omega_3(u_n))\}
\label{eq:124}
\end{equation}
for $1\le i\le k$.

Denote $D_{3n}=\bigcap_{i=1}^k(D_{ni}\cap G_{ni})$.

Let $\psi_n\in U_{ni}$. We show that then, if the event $A_n\cup D_{3n}$ holds,
\begin{equation}
J_n=\inf_{|t|>u_n}\int l(t-u)Z_n(u)\,du-\int l(x_{ni}-u)Z_n(u)\,du>0,
\label{eq:125}
\end{equation}
which completes the proof.

Since $|I^{1/2}(\theta_0)(t-x_{ni})|>5C_0\delta u_n$ and $\|x_{ni}-u\|<|I^{1/2}(\theta_0)(x_{ni})-u|<2\delta u_n$ for $u\in U_{ni}$, we have
\begin{align}
J_n
&\ge l_1(3\delta u_n)\int_{U_{ni}}Z_n(u)\,du+\inf_{|t|>u_n}\int_{V_{ni}}l(t-u)Z_n(u)\,du\notag\\
&\qquad-l_1(2\delta u_n)\int_{U_{ni}}Z_n(u)\,du-\int_{V_{ni}}l(x_{ni}-u)Z_n(u)\,du\notag\\
&\ge\bigl(l_1(3\delta u_n)-l_1(2\delta u_n)\bigr)\int_{U_{ni}}Z_n(u)\,du-\int_{V_{ni}}l(x_{ni}-u)Z_n(u)\,du\notag\\
&\ge\bigl(l_1(3\delta u_n)-l_1(2\delta u_n)\bigr)\int_{U_{ni}}Z_n(u)\,du-l_1(u_n)\int_{V_{ni}}Z_n(u)\,du.
\label{eq:126}
\end{align}
Using~\eqref{eq:124} and the occurrence of the event $A_n$ we obtain
\begin{equation}
\int_{U_{ni}}Z_n(u)\,du=n^{-d/2}\exp\{|\psi_n|^2/2+o(nu_n^2)\}
\label{eq:127}
\end{equation}
and, since $|u-\psi_n|\ge\delta u_n$ for $u\in V_{ni}$,
\begin{equation}
\int_{V_{ni}}Z_n(u)\,du=o\bigl(n^{-d/2}\exp\{|\psi_n|^2/2-\delta^2 nu_n^2(1+o(1))/2\}\bigr).
\label{eq:128}
\end{equation}
From~\eqref{eq:126}--\eqref{eq:128} relation~\eqref{eq:125} follows.

Thus~\eqref{eq:89} holds if we take $D_n=A_n\cap D_{1n}\cap D_{2n}\cap D_{3n}$.

Consider the case~\eqref{eq:99}. For $\theta\in\Theta$ denote $\lambda=3(b^2+1+d/2)$. Put $\delta_n\asymp n^{-\lambda}$.

With this choice of $\delta_n$, if we prove the analogues of inequalities~\eqref{eq:108},~\eqref{eq:115} and~\eqref{eq:124}, then all the estimates~\eqref{eq:112},~\eqref{eq:119} and~\eqref{eq:126}--\eqref{eq:128} remain valid.

Define the events
\begin{equation}
D_{21n}=D_{2n}(r_1,r_2)=\bigl\{X_1,\dots,X_n:|L_{1n}|>n^{-\lambda/3}\bigr\}
\label{eq:129}
\end{equation}
and put $D_{2n}=D_n(0.5,r_n)$.

For $1\le i\le k$ define the events
\begin{align}
D_{ni}&=\bigl\{X_1,\dots,X_n:|L_{ni}|>n^{-\lambda/3}\bigr\},
\label{eq:130}\\
G_{ni}&=\bigl\{X_1,\dots,X_n:|\overline{L}_{ni}|>n^{-\lambda/3}\bigr\}.
\label{eq:131}
\end{align}
Arguing as in~\eqref{eq:102} we obtain
\begin{equation}
\mathbb{E} L_{1n}<Cnn^{-\lambda}(r_nu_n)^d
\label{eq:132}
\end{equation}
and for $1\le i\le k$
\begin{equation}
\mathbb{E} L_{ni}<Cnn^{-\lambda}u_n^d,\qquad
\mathbb{E}\overline{L}_{ni}<Cnn^{-\lambda}(r_nu_n)^d.
\label{eq:133}
\end{equation}
Hence, estimating as in~\eqref{eq:108},
\begin{equation}
P(D_{21n})<Cn^{1-2\lambda/3}(r_nu_n)^d,\qquad
P(D_{22n})<Cn^{1-2\lambda/3}(r_nu_n)^d
\label{eq:134}
\end{equation}
and for $1\le i\le k$
\begin{equation}
P(D_{2ni})<Cn^{1-2\lambda/3}u_n^d,\qquad
P(G_{2ni})<Cn^{1-2\lambda/3}u_n^d.
\label{eq:135}
\end{equation}
Since these asymptotics are of substantially smaller order than the asymptotics
\begin{equation}
P_{\theta_0}\Biggl(I^{-1/2}(\theta_0)\sum_{i=1}^n\phi_\varepsilon(X_i)>nu_n\Biggr),
\label{eq:136}
\end{equation}
the upper bound in~\eqref{eq:33} follows.

We prove~\eqref{eq:34}--\eqref{eq:36} for $b=0$. In~\eqref{eq:104} we may always take $\delta_0<\delta$. Then, by~\eqref{eq:103},~\eqref{eq:112},~\eqref{eq:113} and~\eqref{eq:119},
\begin{equation}
\lim_{n\to\infty}(nu_n^2/2)^{-1}\log P_{\theta_0}\bigl(\bigl|(\tilde\theta_n-\theta)-2n^{-1/2}\psi_n\bigr|>\delta u_n,\ \Upsilon_n\bigr)=-\infty,
\label{eq:137}
\end{equation}
\begin{equation}
\lim_{n\to\infty}(nu_n^2/2)^{-1}\log P_{\theta_0}\Biggl(\Biggl|n^{-1}\sum_{i=1}^n\xi_n(X_i,\tilde\theta_n)-\psi_n^T\psi_n/2\Biggr|>\delta nu_n^2,\ \Upsilon_n\Biggr)=-\infty
\label{eq:138}
\end{equation}
and
\begin{equation}
\lim_{n\to\infty}(nu_n^2/2)^{-1}\log P_{\theta_0}\Biggl(\Biggl|\sum_{i=1}^n\xi_n(X_i,\tilde\theta_n)
-n(\tilde\theta_n-\theta_0)^TI(\theta)(\tilde\theta_n-\theta_0)/2\Biggr|>\delta nu_n^2,\ \Upsilon_n\Biggr)=-\infty.
\label{eq:139}
\end{equation}
If the complementary event $\overline{\Upsilon}_n$ holds, inequalities analogous to~\eqref{eq:137}--\eqref{eq:139} follow from the estimates~\eqref{eq:126}--\eqref{eq:128}, which completes the proof of~\eqref{eq:34}--\eqref{eq:36}.

From the estimates~\eqref{eq:103},~\eqref{eq:112},~\eqref{eq:113},~\eqref{eq:119} and~\eqref{eq:126}--\eqref{eq:128} we obtain that there exists $c>0$ such that
\begin{equation}
\lim_{n\to\infty}(nu_n^2)^{-1}\log P_\theta\Biggl(\int_{|u|<r_nu_n}l_1(\|\tilde\theta_n-u\|)Z_n(u)\,du\le\exp\{-cnu_n^2\}\Biggr)=-\infty.
\label{eq:140}
\end{equation}
Consequently, in the proof of the theorem we may neglect the contribution of the terms $I(r_n)$, $J(r_n)$, $Q(r_n)$ when performing the estimates.

\subsection{Proof of Theorem~\ref{th:posterior}}

The arguments are largely based on the estimates appearing in the proof of Theorem~\ref{th:Bayes}. We shall adhere to the notation of the proof of Theorem~\ref{th:Bayes}.

We begin with the proof of the lower bound in~\eqref{eq:42}. Denote $r=\inf\{|\tau|:\tau\in\Omega\}$. For any $\varepsilon>0$ there exist a point $e\in\Omega$ and $\delta>0$ such that $|e|-r<\varepsilon$ and $B(e,2\delta)\subset\Omega$.

Define the events
\begin{equation}
\Psi_{ne}=\{X_1,\dots,X_n:\psi_n\in n^{1/2}u_n B(e,\delta)\}.
\label{eq:141}
\end{equation}
Then by Theorem~\ref{th:LD}
\begin{equation}
\lim_{n\to\infty}(nu_n^2/2)^{-1}\log P(\Psi_{ne})=-\inf_{x\in B(e,\delta)}|x|^2.
\label{eq:142}
\end{equation}
Define the event
\begin{equation}
G_n=\Biggl\{X_1,\dots,X_n:(\psi_n^T\psi_n/2)^{-1}\log\int_{\overline{B}(une,5\delta u_n)}Z_n(u)\pi(u)\,du>1-\delta\Biggr\}.
\end{equation}
Then, estimating analogously to~\eqref{eq:124},~\eqref{eq:127} and~\eqref{eq:128}, we obtain
\begin{equation}
\lim_{n\to\infty}(nu_n^2/2)^{-1}\log P(\Psi_{ne}\cap G_n)=-\infty.
\label{eq:143}
\end{equation}
For $0<\varepsilon<1$ define the event
\begin{equation}
G_{1n\varepsilon}=\Biggl\{X_1,\dots,X_n:\Bigl|(\psi_n^T\psi_n/2)^{-1}\log\int_{B(une,\delta u_n)}Z_n(u)\pi(u)\,du-1\Bigr|<\varepsilon\Biggr\}.
\end{equation}
Then, estimating analogously to~\eqref{eq:124},~\eqref{eq:127} and~\eqref{eq:128}, we obtain that for every $\varepsilon>0$
\begin{equation}
\lim_{n\to\infty}(nu_n^2/2)^{-1}\log P(\Psi_{ne}\cap G_{1n\varepsilon})=-\infty.
\label{eq:144}
\end{equation}
From~\eqref{eq:142}--\eqref{eq:144} the lower bound in Theorem~\ref{th:posterior} follows.

Applying the methods of the proof of Theorem~\ref{th:Bayes}, namely the inequalities~\eqref{eq:108},~\eqref{eq:111},~\eqref{eq:115},~\eqref{eq:118} when the event $\Upsilon_n$ holds, and the equalities~\eqref{eq:127},~\eqref{eq:128} when the event $\overline{\Upsilon}_n$ holds, we obtain that for every $C>0$ and for the set
\begin{equation}
G_{2n}=\{X_1,\dots,X_n:\psi_n<(r-\delta)n^{1/2}u_n\}
\label{eq:145}
\end{equation}
one has
\begin{equation}
\lim_{n\to\infty}(nu_n^2/2)^{-1}\log P\Biggl(\frac{\displaystyle\int_{|u|<ru_n}Z_n(u)\pi(u)\,du}{\displaystyle\int_{|u|>ru_n}Z_n(u)\pi(u)\,du}<C,\ G_{2n}\Biggr)=-\infty,
\label{eq:146}
\end{equation}
which, together with Theorem~\ref{th:LD}, yields the upper bound in Theorem~\ref{th:posterior}.

\begin{thebibliography}{99}

\bibitem{1}
M.\,A.~Arcones.
Moderate deviations for $M$-estimators.
\textit{Test}, 11:465--500, 2002.

\bibitem{2}
A.\,A.~Borovkov and A.\,A.~Mogulskii.
Large deviations and testing statistical hypotheses.
\textit{Siberian Advances in Mathematics}, 2 (1992); 3 (1993).

\bibitem{3}
M.\,S.~Ermakov.
Asymptotically efficient statistical inference for moderate deviation probabilities.
\textit{Theory of Probability and its Applications}, 48:622--641, 2004.

\bibitem{4}
M.\,S.~Ermakov.
The sharp lower bounds of asymptotic efficiency of estimators in the zone of moderate deviation probabilities.
\textit{Electronic Journal of Statistics}, 6:2150--2184, 2012.

\bibitem{5}
M.\,S.~Ermakov.
Bahadur asymptotic efficiency in the zone of moderate deviation probabilities.
arXiv:2504.19331, 14 pages, 2025.

\bibitem{6}
M.\,S.~Ermakov.
On the lower bound of the exact asymptotics of probabilities of large deviations of statistical estimators.
\textit{Problems of Information Transmission}, 35(3):54--66, 1999.

\bibitem{7}
M.\,S.~Ermakov.
Local asymptotic normality of the logarithm of the likelihood ratio in the zone of moderate deviation probabilities.
\textit{Zapiski Nauchnykh Seminarov POMI}, 525:71--85, 2023.

\bibitem{8}
F.\,Q.~Gao.
Moderate deviations for the maximum likelihood estimator.
\textit{Statistics \& Probability Letters}, 55:345--352, 2001.

\bibitem{9}
J.~Hajek.
Local asymptotic minimax and admissibility in estimation.
In \textit{Proceedings of the Sixth Berkeley Symposium on Mathematical Statistics and Probability}, volume~1, pages 175--194, 1972.

\bibitem{10}
I.\,A.~Ibragimov and R.\,Z.~Hasminskii.
\textit{Statistical Estimation: Asymptotic Theory}.
Springer, New York, 1981.

\bibitem{11}
I.\,A.~Ibragimov and M.~Radavicius.
Probability of large deviations for the maximum likelihood estimator.
\textit{Soviet Mathematics Doklady}, 23(2):403--406, 1981.

\bibitem{12}
R.\,W.~Keener.
\textit{Theoretical Statistics}.
Springer, Berlin, 2010.

\bibitem{13}
L.~Le Cam.
\textit{Asymptotic Methods in Statistical Decision Theory}.
Springer, Berlin, 1986.

\bibitem{14}
Y.~Miao and Y.\,X.~Chen.
Note on the moderate deviation principle of maximum likelihood estimator.
\textit{Acta Applicandae Mathematicae}, 110(2):863--869, 2010.

\bibitem{15}
Yu.\,V.~Prokhorov.
An extremal problem in probability theory.
\textit{Theory of Probability and its Applications}, 4:211--214, 1959.

\bibitem{16}
I.\,M.~Skovgaard.
Large deviation approximations for maximum likelihood estimators.
\textit{Probability and Mathematical Statistics}, 6:89--107, 1985.

\bibitem{17}
Z.\,H.~Xiao and L.\,Q.~Liu.
Moderate deviations of maximum likelihood estimator for independent not identically distributed case.
\textit{Statistics \& Probability Letters}, 76:1056--1064, 2006.

\bibitem{18}
A.\,W.~van der Vaart.
\textit{Asymptotic Statistics}.
Cambridge University Press, Cambridge, 1998.

\end{thebibliography}

\end{document}



